% 英文摘要
\textbf{Power Semiconductor Industry Deep Dive: AI + New Energy Dual Engine, Third-Generation Semiconductors Enter Accelerated Penetration Phase}

\vspace{0.5em}

The global power semiconductor industry stands at a historic inflection point, driven by two secular growth engines: AI-driven data center power upgrade and the global energy transition. We estimate the global power semiconductor market will reach approximately \$54–56 billion in 2025, with a 5–6\% CAGR through 2030, while the China market will grow at a faster 10–12\% CAGR to reach \textyen190–200 billion in 2025, supported by both volume growth and domestic substitution.

\vspace{0.3em}

\textbf{AI Data Centers: The Most Significant Marginal Growth Driver.} As AI server racks evolve from traditional CPU-based systems (~\$2,000–5,000 per rack in power semiconductor value) to Hopper/H100 GPU clusters (\$30,000–40,000 per rack) and next-generation Rubin/Blackwell architectures (\$50,000–80,000 per rack), power semiconductor content per rack increases by 10–15×. Key technology trends include the shift to 48V bus architectures (expected to reach ~80\% penetration by 2027), GaN-based VRM adoption (projected 22–30\% penetration by 2026), and overall system voltage increases. We forecast AI server power semiconductor revenue will grow from ~\$3.9 billion in 2024 to ~\$9.8 billion in 2029, representing a ~20\% CAGR.

\vspace{0.3em}

\textbf{New Energy Vehicles: The Largest End-Market.} EVs remain the single largest downstream application for power semiconductors, with silicon IGBTs as the mainstream solution and SiC MOSFETs gaining rapid share in 800V high-voltage platforms. We expect automotive-grade IGBT and SiC demand to continue growing at a 15–20\% CAGR as EV penetration rises and per-vehicle power semiconductor content increases.

\vspace{0.3em}

\textbf{Third-Generation Semiconductors: Accelerated Penetration.} SiC and GaN are no longer niche technologies. We project the global SiC power device market will reach \$3.0–3.5 billion in 2025 and \$8.5–10 billion by 2030 (23–25\% CAGR), while GaN power devices will grow from \$0.8–0.9 billion in 2025 to \$3.5–4.3 billion by 2030 (35–40\% CAGR). These technologies are complementary rather than substitutive—SiC dominates high-voltage/high-power applications (EV traction inverters, renewables) while GaN leads in high-frequency scenarios (data center VRM, fast chargers, consumer electronics).

\vspace{0.3em}

\textbf{Domestic Substitution: From Consumer to Automotive-Grade.} China's overall power semiconductor self-sufficiency rate stands at approximately 30–35\%, with wide variation across segments: 50–70\% for low-voltage MOSFETs, but only 10–15\% for automotive-grade IGBTs. The progression from consumer/industrial to automotive-grade represents the most critical value-upgrade path for domestic players. Companies with IDM capabilities, automotive qualification, and third-generation semiconductor R\&D pipelines are best positioned.

\vspace{0.3em}

\textbf{Investment Thesis: Overweight.} We initiate coverage of China's power semiconductor sector with an \textbf{Overweight} rating. We recommend a three-pronged investment strategy: (1) IDM leaders with strong profitability and manufacturing scale; (2) automotive-grade IGBT/SiC leaders with customer validation; and (3) high-growth third-generation semiconductor substrate and device players. Our top picks include CRRC Times Electric (\textbf{Buy}), StarPower (\textbf{Overweight}), and CR Micro (\textbf{Overweight}).

\vspace{0.3em}

\textbf{Key Risks:} AI capex slowdown risk, EV demand below expectations, SiC technology iteration slower than expected, intensified price competition, trade policy uncertainties, and geopolitical supply chain disruptions.
